\documentclass[11pt]{article}

\usepackage[a4paper, margin=1in, headheight=13.6pt]{geometry}
\usepackage[T1]{fontenc}
\usepackage{lmodern}
\usepackage{amsmath}
\usepackage{amssymb}
\usepackage{booktabs}
\usepackage{tabularx}
\usepackage{longtable}
\usepackage{enumitem}
\usepackage{fancyhdr}
\usepackage{graphicx}
\usepackage{xcolor}
\usepackage[hidelinks]{hyperref}
\usepackage{pdflscape}
\usepackage[backend=biber,style=numeric-comp,sorting=none,url=true,isbn=false]{biblatex}
\addbibresource{references.bib}

\newcommand{\pt}{p_{\mathrm{T}}}
\newcommand{\GeV}{\ensuremath{\,\mathrm{GeV}}}
\newcommand{\ptll}{\pt^{\ell\ell}}
\newcommand{\ptj}{\pt^{\mathrm{jet}}}
\newcommand{\code}[1]{\texttt{#1}}

\pagestyle{fancy}
\fancyhf{}
\lhead{CMS Analysis Note (draft, not for circulation)}
\rhead{boostedV: \code{ZX\_test} selection}
\cfoot{\thepage}

\title{Energy correlators of boosted jets recoiling against a $Z\to\ell\ell$ boson:\\
event selection, object definitions and corrections of the \code{ZX\_test} configuration}
\author{Simon Rothman}
\date{Draft of \today}

\begin{document}
\sloppy
\maketitle

\begin{abstract}
This note describes, step by step, the \code{ZX\_test} configuration of the boostedV workflow: the simulated samples,
the reconstruction-level and generator-level event selections, the muon momentum assignment and corrections, the AK8
jet processing (jet energy corrections re-applied from the raw momentum, jet energy resolution smearing, PUPPI-weighted
constituents, jet identification, lepton--jet overlap flags and the back-to-back requirement), the energy correlator
observables and their output format, the order in which corrections and selections are applied, the event weights
stored, a cutflow, and the regions of the columnar analysis. Each correction and identification is related to the
corresponding CMS physics object group (POG) recommendation, and deviations from it are stated. It
describes the simulation code at boostedV commit \code{a2ad376} (EventSelection \code{4cbe057}, JetSetup
\code{04df502}), CMSSW\_15\_0\_22, and the columnar analysis as of 15 November 2026, applied to Run~3 2022 simulation
with post-EE conditions. Values marked \emph{placeholder} are development choices, not CMS recommendations. Open points
are collected in Section~\ref{sec:issues}.
\end{abstract}

\tableofcontents
\newpage

%=======================================================================================================================
\section{Introduction and scope}
\label{sec:intro}

The boostedV project aims to measure energy correlators of boosted, hadronically decaying electroweak
bosons~\cite{Ricci:2022htc} and of the quark and gluon jets that mimic them. The \code{ZX\_test} configuration, the first one used for physics studies, selects
events with a $Z\to\mu\mu$ or $Z\to ee$ candidate at high transverse momentum recoiling against an
anti-$k_\mathrm{T}$~\cite{Cacciari:2008gp,Cacciari:2011ma} jet of distance parameter $R = 0.8$ (AK8 jet), clustered
from particle-flow candidates~\cite{CMS:2017yfk} weighted with PUPPI~\cite{Bertolini:2014bba,CMS:2020ebo}. In $ZZ\to\ell\ell q\bar q$ and $WZ\to q\bar q'\ell\ell$ events this
jet is frequently the hadronically decaying boson; in $Z$+jets events it is a quark or gluon jet. For each selected jet,
the configuration stores the two-point projected energy correlator (E2C)~\cite{Chen:2020vvp} with energy weights of power
$n = 1$ and $n = 2$, and the one-point correlator (E1C) in three reference frames, all computed from the charged jet constituents
and stored without binning.

The reconstruction-level selection is deliberately loose: leptons are required to pass kinematic criteria only, without
identification or isolation, so that these requirements and the trigger can be studied offline from the stored lepton
properties (Section~\ref{sec:objects}). Every selection step is evaluated both at reconstruction level (\emph{reco})
and at generator level (\emph{gen}) with the same channel definitions, and an event is retained if it passes at either
level.

The entire configuration is defined in one YAML file, \code{simulation/configs/ZX\_test.yaml}, which is validated
against a JSON schema (\code{analysis.schema.yaml}) and additional consistency checks (\code{analysisConfig.py});
unknown or missing keys are errors. The CMSSW configuration \code{eecNano\_cfg.py} builds the full process from this
file; only per-job settings (input and output files, number of events, data or simulation, and whether constituent
tables are written) are command-line arguments. Configuration keys, module labels and output columns are written in
\code{typewriter} font. The labels \code{<c>} and \code{<ch>} denote a channel name (\code{Zmumu} or \code{Zee}).

%=======================================================================================================================
\section{Software and workflow}
\label{sec:workflow}

\begin{description}[style=nextline]
\item[Release and packages] CMSSW\_15\_0\_22 (\code{el9\_amd64\_gcc12}) with four boostedV packages:
  \code{SRothman/EECs} (correlator library and its EDM producers), \code{SRothman/DataFormats} (the \code{EECJet}
  jet format), \code{SRothman/JetSetup} (jet processing) and \code{SRothman/EventSelection} (event selection, muon
  corrections and event-level tables).
\item[Job] \code{cmsRun eecNano\_cfg.py config=ZX\_test.yaml globalTag=<GT> inputFiles=<MiniAOD> isMC=1
  writeConstituents=1} reads MiniAOD and writes a file in NanoAOD format~\cite{xpogdocs} (\emph{NANO}) and a histogram file. All POG
  corrections are read with correctionlib~\cite{correctionlib} from the central payloads~\cite{jsonpog} under
  \code{/cvmfs/cms-griddata.cern.ch/cat/metadata/}. With
  \code{ZX\_test} the histogram file is empty, because the correlators are stored unbinned and are not aggregated
  (\code{aggregation.observables: []}). The processing chain is summarised in Table~\ref{tab:chain}.
\item[Production] \code{simulation/production}: one CRAB task per sample (\code{submit.py}), with file-based
  splitting, a 3000\,MB memory request and output to \code{T3\_US\_FNALLPC}; each job writes \code{nano\_<i>.root} and
  \code{hists\_<i>.root} and is complete only when both exist. \code{status.py} reports and resubmits failed jobs, and
  \code{merge.py} writes for each sample the list of NANO files of the complete jobs and the sums \code{genEventCount},
  \code{genEventSumw} and \code{genEventSumw2} from the Runs trees of exactly those jobs; it refuses to merge with
  missing jobs unless explicitly allowed. The git state of all repositories, the manifest and the analysis YAML are
  recorded at submission.
\item[Analysis] \code{analysis/} (package \code{bvana}): a coffea processor driven by a second YAML file,
  \code{analysis/configs/zxtest.yaml}, which defines regions, cuts, binnings and histograms (Section~\ref{sec:regions}).
\end{description}

\paragraph{Resources.} Measured locally on Drell--Yan simulation with the configuration described here: a start-up time
of 28\,s, then about 5\,ms per event for inclusive events and 42\,ms per event for events with $\ptll$ between 400 and
600\GeV; peak resident memory 1.9\,GB; about 6.5\,kB of output per input event for the high-$\ptll$ events. Each NANO file
also contains a fixed 2.5\,MB \code{ParameterSets} tree.

\begin{table}[htbp]
\centering
\caption{Processing chain in simulation, in dependency order. Module labels are those of \code{eecNano\_cfg.py}; gen
labels carry the prefix \code{gen}. In data, the gen-level steps and the JER smearing are absent, and the JEC uses the
data key.}
\label{tab:chain}
\small
\begin{tabularx}{\textwidth}{l>{\raggedright\arraybackslash}X>{\raggedright\arraybackslash}X}
\toprule
Step & Modules & Section \\
\midrule
Event filters & \code{goodVertexFlag}, \code{metFilterFlag}, \code{<c>Trigger} & \ref{sec:trigger}, \ref{sec:vertex} \\
Muon momentum & \code{correctedMuons} (\code{MuonCorrectionProducer}) & \ref{sec:muoncorr} \\
Lepton objects & \code{looseMuons}, \code{looseElectrons}; gen: \code{genParticleLevel}, \code{genTightMuons},
  \code{genTightElectrons} & \ref{sec:objects}, \ref{sec:gen} \\
Dilepton & \code{<c>Sel} (\code{DileptonSelector}); flag \code{<c>} (\code{FlagCombiner}) & \ref{sec:dilepton} \\
Reco jets & \code{jetsRecoJEC} $\to$ \code{jetsRecoRaw} $\to$ \code{jetsRecoJER} $\to$ \code{jetsRecoConst} $\to$
  \code{jetsReco} & \ref{sec:recojets} \\
Gen jets & \code{jetsGenRaw} $\to$ \code{jetsGenConst} $\to$ \code{jetsGen} & \ref{sec:genjets} \\
Overlap flags & \code{jetsRecoOverlap}, \code{jetsGenOverlap} (\code{EECJetOverlapProducer}) & \ref{sec:veto} \\
Back-to-back & \code{jets<Coll>BackToBack<c>} (\code{BackToBackSelector}) $\to$ \code{jets<Coll>B2B}
  (\code{EECJetMaskSelector}) & \ref{sec:b2b} \\
Event filter & \code{anyChannel}, \code{genAnyChannel} $\to$ \code{anySelectedFilter} & \ref{sec:overview} \\
Correlators & \code{eec<Coll><obs>} and \code{eec<Coll><obs>Table} & \ref{sec:eec} \\
Tables & jet, lepton, event, weight and generator tables & \ref{sec:output} \\
\bottomrule
\end{tabularx}
\end{table}

%=======================================================================================================================
\section{Samples and normalisation}
\label{sec:samples}

All samples are Run3Summer22EE MiniAODv4 simulation (conditions of the 2022 eras E, F and G), processed with the global
tag \code{130X\_mcRun3\_2022\_realistic\_postEE\_v6}. No collision data have been processed with \code{ZX\_test}.
Table~\ref{tab:samples} lists the samples of the next production (manifest
\code{simulation/production/samples\_zxb2b.yaml}, not yet submitted) and those of the first production (tag
\code{20261113\_zxtest}), which used an earlier version of the configuration without the dilepton $\pt$ and back-to-back
requirements.

\begin{table}[htbp]
\centering
\caption{Samples. Event and file counts from DAS. For \code{20261113\_zxtest}, the processed fraction is the
\code{genEventCount} of the merged jobs divided by the DAS event count; the missing jobs are discussed in
Section~\ref{sec:issues}. For the next production, the splitting is given in input files per CRAB job.}
\label{tab:samples}
\small
\begin{tabularx}{\textwidth}{>{\raggedright\arraybackslash}Xrrl}
\toprule
Dataset (\code{/<name>/Run3Summer22EEMiniAODv4-\allowbreak{}130X\_\allowbreak{}...\_\allowbreak{}postEE\_\allowbreak{}v6-\allowbreak{}v*/MINIAODSIM}) & Events & Files & Jobs \\
\midrule
\multicolumn{4}{l}{\emph{Next production, manifest \code{samples\_\allowbreak{}zxb2b.yaml} (configuration of this note)}} \\
\code{ZZto2L2Q\_\allowbreak{}TuneCP5\_\allowbreak{}13p6TeV\_\allowbreak{}powheg-\allowbreak{}pythia8} & 49\,084\,483 & 891 & 6 files/job \\
\code{WZto2L2Q\_\allowbreak{}TuneCP5\_\allowbreak{}13p6TeV\_\allowbreak{}powheg-\allowbreak{}pythia8} & 14\,954\,528 & 278 & 6 files/job \\
\code{DYto2L-\allowbreak{}4Jets\_\allowbreak{}MLL-\allowbreak{}50\_\allowbreak{}PTLL-\allowbreak{}200to400\_\allowbreak{}TuneCP5\_\allowbreak{}13p6TeV\_\allowbreak{}madgraphMLM-\allowbreak{}pythia8} & 6\,439\,323 & 286 & 4 files/job \\
\code{DYto2L-\allowbreak{}4Jets\_\allowbreak{}MLL-\allowbreak{}50\_\allowbreak{}PTLL-\allowbreak{}400to600\_\allowbreak{}TuneCP5\_\allowbreak{}13p6TeV\_\allowbreak{}madgraphMLM-\allowbreak{}pythia8} & 3\,669\,178 & 180 & 4 files/job \\
\code{DYto2L-\allowbreak{}4Jets\_\allowbreak{}MLL-\allowbreak{}50\_\allowbreak{}PTLL-\allowbreak{}600\_\allowbreak{}TuneCP5\_\allowbreak{}13p6TeV\_\allowbreak{}madgraphMLM-\allowbreak{}pythia8} & 3\,234\,438 & 186 & 4 files/job \\
\midrule
\multicolumn{4}{l}{\emph{First production \code{20261113\_\allowbreak{}zxtest} (earlier configuration), merged}} \\
\code{DYto2L-\allowbreak{}4Jets\_\allowbreak{}MLL-\allowbreak{}50\_\allowbreak{}TuneCP5\_\allowbreak{}13p6TeV\_\allowbreak{}madgraphMLM-\allowbreak{}pythia8} & 264.4\,M & 6114 & 589/613, 96.0\% \\
\code{TTto2L2Nu\_\allowbreak{}TuneCP5\_\allowbreak{}13p6TeV\_\allowbreak{}powheg-\allowbreak{}pythia8} & 74.3\,M & 1388 & 229/232, 98.7\% \\
\code{WZ\_\allowbreak{}TuneCP5\_\allowbreak{}13p6TeV\_\allowbreak{}pythia8} & 26.7\,M & 391 & 79/79, 100\% \\
\code{ZZ\_\allowbreak{}TuneCP5\_\allowbreak{}13p6TeV\_\allowbreak{}pythia8} & 4.04\,M & 82 & 17/17, 100\% \\
\bottomrule
\end{tabularx}
\end{table}

\paragraph{Choice of samples.} The project plan requires $Z(\ell\ell)$+jets, $WZ$ and $ZZ$ samples, the latter two
with $Z\to\ell\ell$ and a hadronically decaying $W$ or $Z$ boson. In the first production, the inclusive $ZZ$ sample
yields only 960 events in the boosted region of Section~\ref{sec:regions}. The $ZZ\to2\ell2q$ sample would correspond
to about 130 times the equivalent inclusive statistics if ``2L'' denoted electrons and muons only; however, the reco
channel flags are passed by only 7.5\% ($\mu\mu$) and 5.6\% ($ee$) of its events (Table~\ref{tab:cutflow}), which
suggests that $Z\to\tau\tau$ decays are included and the gain is smaller. This is to be checked in the generator
configuration. The Drell--Yan samples binned in \code{PTLL} provide 13.3\,M events in the region $\ptll > 150\GeV$
selected here (Section~\ref{sec:dilepton}), compared with about 0.5\,M events with $\ptll > 200\GeV$ in the inclusive
sample (0.21\% of its selected $Z\to\mu\mu$ events). \code{PTLL} is a generator-level dilepton $\pt$ (to be confirmed
from the generator configuration); combining these samples with the inclusive one requires an overlap removal in that
variable.

\paragraph{Normalisation.} The NANO stores \code{genWeight} per event, and the Runs tree stores \code{genEventCount},
\code{genEventSumw} and \code{genEventSumw2} summed over all processed events (the weight table is scheduled on its own
path, independent of the event filter). A sample is normalised with
$w = \code{genWeight}\cdot\sigma\cdot L/\code{genEventSumw}$, where \code{genEventSumw} is summed over the merged jobs
only, so that missing jobs do not bias the normalisation. Cross sections have not yet been collected
(\code{analysis/xsec.yaml}); plots produced so far are normalised to unit area.

%=======================================================================================================================
\section{Event selection: overview}
\label{sec:overview}

\code{ZX\_test} defines two channels, \code{Zmumu} and \code{Zee}, both of type \code{dilepton}, at two levels:
\begin{itemize}
\item \textbf{reco}: built from MiniAOD objects; a channel requires its trigger, a good primary vertex, the MET filters
  and the dilepton selection;
\item \textbf{gen} (simulation only): built from Rivet particle-level objects; a channel requires the dilepton
  selection only.
\end{itemize}
Each step is a boolean EDM product; Table~\ref{tab:flags} lists the flags, which are stored in the singleton tables
\code{Selection} (reco) and \code{GenSelection} (gen). An event is written if \code{B2B<c>} is true for at least one
channel at one level (\code{selection.levels: [reco, gen]}; module \code{anySelectedFilter}). Because both levels are
evaluated for every event in simulation, every written event carries the flags and selected objects of both levels, so
that events migrating across the selection between the levels are retained. The order in which the corrections and
selections act on each object is summarised in Section~\ref{sec:order}.

\begin{table}[htbp]
\centering
\caption{Selection flags per channel \code{<c>}. Columns \code{Selection\_<flag>} and \code{GenSelection\_<flag>}.}
\label{tab:flags}
\small
\begin{tabularx}{\textwidth}{lc>{\raggedright\arraybackslash}X}
\toprule
Flag & Levels & Definition \\
\midrule
\code{goodVertex} & reco & first primary vertex good (Section~\ref{sec:vertex}) \\
\code{metFilters} & reco & all MET filters pass (Section~\ref{sec:vertex}) \\
\code{<c>Trigger} & reco & any of the channel's HLT paths fired (Section~\ref{sec:trigger}) \\
\code{<c>Sel} & both & dilepton selection, including $\ptll \geq 150\GeV$ (Section~\ref{sec:dilepton}) \\
\code{<c>} & both & reco: \code{<c>Sel} AND \code{<c>Trigger} AND \code{goodVertex} AND \code{metFilters};
  gen: \code{<c>Sel} \\
\code{B2B<c>} & both & \code{<c>} AND a back-to-back jet at the same level (Section~\ref{sec:b2b}) \\
\bottomrule
\end{tabularx}
\end{table}

%=======================================================================================================================
\section{Reconstruction-level event selection}
\label{sec:reco}

\subsection{Trigger}
\label{sec:trigger}
The trigger decision is read from the \code{TriggerResults} of process \code{HLT}. Paths are specified as glob
patterns with the version suffix \code{\_v*}; \code{<c>Trigger} is true if any matching path fired. A pattern that
matches no path of the menu is a configuration error (\code{allowMissingTriggers: false}). The paths are
\begin{itemize}
\item \code{Zmumu}: \code{HLT\_IsoMu24}, \code{HLT\_Mu50}, \code{HLT\_CascadeMu100}, \code{HLT\_HighPtTkMu100};
\item \code{Zee}: \code{HLT\_Ele32\_WPTight\_Gsf}, \code{HLT\_Ele115\_CaloIdVT\_GsfTrkIdT}, \code{HLT\_Photon200}.
\end{itemize}
The selection does not require trigger matching. Each selected lepton carries the column \code{triggerMatched}, true if
the lepton lies within $\Delta R < 0.1$ of an HLT object from the last filter of a fired path among the channel's
paths. No trigger efficiency corrections are applied (Section~\ref{sec:weights}).

\subsection{Primary vertex and MET filters}
\label{sec:vertex}
\begin{itemize}
\item \code{goodVertex}: the first vertex of \code{offlineSlimmedPrimaryVertices} exists, is not fake, and has
  $\mathrm{ndof} > 4$, $|z| < 24\,\mathrm{cm}$ and $\rho < 2\,\mathrm{cm}$.
\item \code{metFilters}: the Run~3 MET-filter list of the JetMET POG~\cite{JMEmetfilters}, i.e.\ all of
  \code{Flag\_goodVertices}, \code{Flag\_globalSuperTightHalo2016Filter},
  \code{Flag\_EcalDeadCellTriggerPrimitiveFilter}, \code{Flag\_BadPFMuonFilter}, \code{Flag\_BadPFMuonDzFilter},
  \code{Flag\_hfNoisyHitsFilter}, \code{Flag\_eeBadScFilter} and \code{Flag\_ecalBadCalibFilter}, read from the
  \code{TriggerResults} of process \code{PAT}, must pass.
\end{itemize}

\subsection{Muon momentum assignment and corrections}
\label{sec:muoncorr}
Muons are reconstructed as described in Ref.~\cite{CMS:2018rym}. Before any muon selection, each \code{slimmedMuons}
muon is copied by \code{MuonCorrectionProducer} (\code{muonCorrections.apply: true}); all muon consumers (object
selection, overlap flags and tables) read these copies. The procedure follows the MUO POG Run~3
recommendations~\cite{MUOcorrections,MUOscare,MUOrecommendations}, which distinguish a medium-$\pt$ regime (scale and
resolution corrections) and a high-$\pt$ regime (TuneP momentum assignment and generalised-endpoint corrections); for
analyses with muons on both sides of $200\GeV$, they recommend correcting each muon with the method of its
regime~\cite{MUOcorrections}.
\begin{enumerate}
\item \textbf{Momentum assignment} (\code{momentum: hybrid}, \code{regimeThreshold: 200}): a muon with
  uncorrected PF $\pt \geq 200\GeV$ takes the TuneP transverse momentum~\cite{CMS:2019ied,MUOrecommendations} and is assigned to the \emph{high-$\pt$ regime}
  (\code{highPtRegime} = 1); otherwise it keeps the PF momentum and is assigned to the \emph{medium regime} (0). In both
  regimes the direction ($\eta$, $\phi$) of the PF muon is kept; only $\pt$ changes. The input $\pt$ is stored as
  \code{originalPt}.
\item \textbf{Medium regime}: the MUO POG Run~3 ``MuonScaRe'' scale correction (Z-peak method) and, in simulation, the
  resolution smearing, from \code{MUO/\allowbreak{}Run3-\allowbreak{}22EFGSep23-\allowbreak{}Summer22EE-\allowbreak{}NanoAODv12/\allowbreak{}latest/\allowbreak{}muon\_scalesmearing.json.gz}, for
  $26 \leq \pt \leq 200\GeV$ (\code{lowPtThreshold}, \code{highPtThreshold}; evaluated on the $\pt$ before
  correction), the range in which the corrections were derived~\cite{MUOscare}; outside it the momentum is unchanged.
  The scale correction is applied first, then the smearing, as in the POG example. The corrections were derived for
  tightly identified and isolated muons~\cite{MUOscare}, whereas here they are applied to all muons, before any
  identification or isolation requirement. The implementation is a port of the POG example \code{MuonScaRe.cc}; as in the POG code, the smearing
  random number depends on the event number only, so all muons of an event share it.
\item \textbf{High regime} (uncorrected TuneP $\pt \geq 200\GeV$ and $|\eta| < 2.4$; \code{highPt.era: 2022postEE}):
  the generalised-endpoint curvature bias $\kappa_b$ from the MUO POG 2022 post-EE tables~\cite{MUOcorrections}, which
  for Run~3 are recommended as a correction, applied to simulation as
  $q/\pt \to q/\pt + \kappa_b$ (\code{kappaTarget: mc}), followed by the simulation smearing
  $p \to p\,(1 + \mathcal{N}(0, f\,\sigma(p,\eta)))$ with the POG resolution parametrisation
  $\sigma(p,\eta)$ and $f = 0.32$ for $|\eta| < 1.2$ and $0.46$ otherwise (simulation about 5\% and 10\% better than data,
  respectively, for 2022 post-EE)~\cite{MUOcorrections}, with the random number seeded per run, luminosity block, event
  and muon index.
\end{enumerate}
The nominal corrections are applied; no systematic variation is evaluated (\code{variation: nominal}). The formula and
sign of the $\kappa_b$ correction, and the smearing and its systematic variations in all regimes, are provisional
pending clarification with the MUO POG (Section~\ref{sec:issues}).

\subsection{Electrons}
Electrons are reconstructed as described in Ref.~\cite{CMS:2020uim}. \code{slimmedElectrons} are used with their
MiniAOD four-momentum, which includes the energy regression of the reconstruction. The EGamma POG recommends in
addition, for Run~3, an energy scale correction in data and a resolution smearing in simulation~\cite{EGMsas}; these are
\emph{not} applied in \code{ZX\_test} (Section~\ref{sec:issues}).

\subsection{Lepton objects}
\label{sec:objects}
Leptons are not dressed with photons at reco level (\code{dressing.mode: none}). The object selections are kinematic
only:
\begin{itemize}
\item \code{looseMuons}: corrected muons with $\pt > 15\GeV$ and $|\eta| < 2.4$;
\item \code{looseElectrons}: $\pt > 15\GeV$ and supercluster pseudorapidity $|\eta_\mathrm{SC}| < 2.5$, excluding the
  barrel--endcap transition $1.4442 < |\eta_\mathrm{SC}| < 1.566$.
\end{itemize}
No identification, isolation or impact-parameter requirement is applied. For offline studies, the selected leptons carry
the properties of Table~\ref{tab:leptoncols}, taken from the PAT object matched to each selected lepton (the nearest
object of the same charge within $\Delta R < 0.1$).

\begin{table}[htbp]
\centering
\caption{Properties stored for the selected reconstructed leptons (tables \code{<c>Leptons}) in addition to $\pt$,
$\eta$, $\phi$, mass, charge, PDG identifier, \code{undressedPt} and \code{triggerMatched}.}
\label{tab:leptoncols}
\small
\begin{tabularx}{\textwidth}{l>{\raggedright\arraybackslash}X}
\toprule
Muons~\cite{MUOrecommendations} & \code{dxy}, \code{dz} (with respect to the first primary vertex); \code{highPtRegime}; \code{tunePPt}
  (uncorrected TuneP $\pt$); \code{originalPt}; \code{pfRelIso04} ($\Delta\beta$-corrected PF isolation,
  $\Delta R = 0.4$); \code{tkRelIso} (tracker isolation, $\Delta R = 0.3$, relative to the TuneP $\pt$);
  \code{looseId}, \code{mediumId}, \code{tightId}, \code{highPtId}, \code{trkHighPtId}; \code{pfIsoId} (VeryLoose to
  VeryVeryTight as 1--6); \code{tkIsoId} (Loose 1, Tight 2) \\
Electrons & \code{dxy}, \code{dz}; \code{scEta}; \code{cutBased} (RunIIIWinter22 V1~\cite{EGMcutid}: 0 fail, 1 veto,
  2 loose, 3 medium, 4 tight); \code{pfRelIso03} ($\Delta\beta$-corrected, $\Delta R = 0.3$); \code{r9}; \code{heepId}
  (HEEP V7.1); \code{mvaIsoWP80},
  \code{mvaIsoWP90}, \code{mvaNoIsoWP80}, \code{mvaNoIsoWP90} (RunIIIWinter22 V1~\cite{EGMmvaid}; the no-isolation
  decisions stored in Run~3 MiniAOD include isolation in the barrel~\cite{EGMmvaid}) \\
\bottomrule
\end{tabularx}
\end{table}

\subsection{Dilepton selection}
\label{sec:dilepton}
All pairs of selected leptons of the channel's collection are considered. A pair is accepted if
\begin{itemize}
\item the leading lepton has $\pt \geq 26\GeV$ (\code{Zmumu}) or $\pt \geq 35\GeV$ (\code{Zee}), and the subleading
  lepton $\pt \geq 15\GeV$;
\item the leptons have opposite charges;
\item the invariant mass satisfies $61.19 \leq m_{\ell\ell} \leq 121.19\GeV$, i.e.\ $m_Z \pm 30\GeV$ with
  $m_Z = 91.1876\GeV$.
\end{itemize}
Among the accepted pairs, the one with mass closest to $m_Z$ is chosen. The channel then requires this pair to have
$\ptll \geq 150\GeV$ (\code{pairPtMin}, placeholder); if it does not, the channel fails, even if another accepted pair
has a larger $\pt$. The $Z$ candidate (table \code{<c>Boson}) is the sum of the two lepton four-momenta, and the two
leptons (table \code{<c>Leptons}) are stored with the leading one first.

The $\ptll$ requirement strongly reduces the number of written events. In the first production, which did not apply it,
the generator-weighted fraction of selected $Z\to\mu\mu$ events with $\ptll > 150\GeV$ was 0.5\% for inclusive
Drell--Yan, 2.5\% for $t\bar t\to 2\ell2\nu$ and about 5\% for the inclusive $WZ$ and $ZZ$ samples.

%=======================================================================================================================
\section{Generator-level event selection}
\label{sec:gen}

The generator-level objects are those of the Rivet~\cite{Bierlich:2019rhm} particle level, computed as in
NanoAOD~\cite{xpogdocs}: \code{ParticleLevelProducer}
runs on the merged \code{prunedGenParticles} and \code{packedGenParticles}, converted to HepMC. It provides
\begin{itemize}
\item prompt final-state leptons (\code{usePromptFinalStates: true}), dressed with all photons within $\Delta R < 0.1$
  (\code{lepConeSize}), with $\pt > 1\GeV$ and $|\eta| < 2.5$;
\item isolated prompt photons ($\pt > 1\GeV$, $|\eta| < 2.5$, relative isolation below 0.5 within $\Delta R = 0.4$) and
  the missing transverse momentum of prompt neutrinos, neither of which is used by \code{ZX\_test};
\end{itemize}
from input particles with $|\eta| < 5$ and no $\pt$ threshold. The gen objects mirror the reco kinematic requirements:
\code{genTightMuons} ($|\mathrm{PDG\,ID}| = 13$, $\pt > 15\GeV$, $|\eta| < 2.4$) and \code{genTightElectrons}
($|\mathrm{PDG\,ID}| = 11$, $\pt > 15\GeV$, $|\eta| < 2.5$ outside $1.4442 < |\eta| < 1.566$). The gen channels use the
same dilepton selection as reco (lepton $\pt$ thresholds, mass window, choice of the pair closest to $m_Z$ and
$\ptll \geq 150\GeV$), without any trigger, vertex or MET-filter requirement.

%=======================================================================================================================
\section{Jets}
\label{sec:jets}

One reco and one gen AK8 jet collection are built, both from MiniAOD jets with the jet daughters as constituents. Jets
keep the order of the MiniAOD collection, and the index in that collection is stored as \code{srcIdx}. The processed
jets are stored in the \code{EECJet} format (four-momentum, raw $\pt$, \code{srcIdx}, jet identification bits, JER
factor and constituents).

\subsection{Reconstructed AK8 jets}
\label{sec:recojets}
The PUPPI AK8 jets \code{slimmedJetsAK8} are processed in the following order, with the corrections recommended by
the JetMET POG for Summer22EE~\cite{JMEjerc}; the methods are described in Ref.~\cite{CMS:2016lmd}.
\begin{enumerate}
\item \textbf{JEC re-application} (\code{PATJetJECProducer}). The raw four-momentum is obtained with
  \code{correctedJet("Uncorrected")}, and the V4 compound correction \code{Summer22EE\_22Sep2023\_V4\_MC\_L1L2L3Res\_AK8PFPuppi}
  from \code{JME/Run3-22EFGSep23-Summer22EE-NanoAODv12/latest/fatJet\_jerc.json.gz} is evaluated with the jet area,
  raw $\pt$, $\eta$ and $\rho$ (\code{fixedGridRhoFastjetAll}); in data, the \code{DATA} key is used, with the run
  number as an additional input. The jet four-momentum is the raw four-momentum times this correction. The MiniAOD
  correction levels of the copy are rescaled consistently, so that \code{rawFactor} and the energy fractions refer to
  the re-corrected jet.
\item \textbf{Translation} (\code{PATJetToEECJetProducer}). Jets with corrected $\pt < 180\GeV$ are discarded. Each
  daughter (a \code{packedPFCandidates} entry) becomes a constituent whose four-momentum is scaled by its PUPPI weight
  (\code{pat::PackedCandidate::puppiWeight()}); constituents with vanishing scaled three-momentum are skipped. A jet
  above the threshold without daughters stops the job with an error, since MiniAOD keeps the daughters of
  \code{slimmedJetsAK8} only for jets above $170\GeV$ with the MiniAOD correction (Section~\ref{sec:issues}). The jet
  identification bits are computed with \code{PFJetIDSelectionFunctor}, version \code{RUN3PUPPIruns2022FGruns2023CD}:
  bit 0 for the TIGHT and bit 1 for the TIGHTLEPVETO criteria of the JetMET POG Run~3 PUPPI jet
  identification~\cite{JMEjetid}. They are computed from the energy fractions and multiplicities, which the JEC does not
  change.
\item \textbf{JER smearing} (simulation only) with the hybrid method recommended by the JetMET POG~\cite{JMEjerc} and
  the keys
  \code{Summer22EE\_22Sep2023\_JRV2\_MC\_ScaleFactor\_AK8PFPuppi} and \code{\dots\_PtResolution\_AK8PFPuppi}. The scale factor $s(\eta, \pt)$ and the resolution
  $\sigma(\eta, \pt, \rho)$ are evaluated at the corrected $\pt$ of the previous step. A jet is matched to the closest
  \code{slimmedGenJetsAK8} jet (all gen-jet constituents) satisfying $\Delta R < R/2 = 0.4$ and
  $|\pt - \pt^\mathrm{gen}| < 3\,\sigma\,\pt$, where $\sigma$ is the relative $\pt$ resolution. A matched jet is scaled
  by $1 + (s-1)(\pt - \pt^\mathrm{gen})/\pt$, an unmatched one by $1 + \mathcal{N}(0,1)\,\sigma\sqrt{\max(s^2-1, 0)}$,
  where $s$ is the resolution scale factor; negative factors are set to zero. Only the jet four-momentum is scaled, not
  the constituents. The random number is seeded per run, luminosity block, event and \code{srcIdx}; the factor is
  stored as \code{jerFactor}.
\item \textbf{Constituent selection}: only charged constituents are kept (\code{chargedOnly: true}; no $\pt$ or PUPPI
  weight threshold).
\item \textbf{Jet selection}: $\pt \geq 200\GeV$ (after JEC and JER), $|\eta| \leq 2.4$, TIGHT jet identification, and
  at least two constituents after the constituent selection, i.e.\ at least two charged constituents.
\end{enumerate}
The JES and JER uncertainty keys (\code{Summer22EE\_22Sep2023\_V4\_MC\_Total\_AK8PFPuppi} and
\code{\dots\_JRV2\_MC\_SFUncertainty\_AK8PFPuppi}) are configured, but no systematic variation is processed
(\code{variations: \{\}}).

\paragraph{Stored jet properties.} The table \code{RecoJet} contains $\pt$, $\eta$, $\phi$ and mass after all
corrections, \code{rawPt}, \code{srcIdx}, \code{jetId} (the two bits), \code{nConstituents} (after the constituent
selection) and \code{jerFactor}. From the re-corrected PAT jet, it also contains \code{jecPt}, \code{jecFactor},
\code{area}, \code{rawFactor}, the $N$-subjettiness variables $\tau_1$--$\tau_4$~\cite{Thaler:2010tr}, the $N_2$ and $N_3$
variables~\cite{Moult:2016cvt} with $\beta = 1$, the PUPPI soft-drop~\cite{Larkoski:2014wba} mass \code{msoftdrop}, the
ParticleNet~\cite{Qu:2019gqs} scores stored in MiniAOD, energy fractions and
multiplicities, the soft-drop subjet indices, and in simulation the hadron and parton flavour and the numbers of $b$ and
$c$ hadrons. The Global Particle Transformer (GloParT~3, based on the Particle Transformer architecture~\cite{Qu:2022mxj})
scores, which are absent from this MiniAOD, are computed in the
job by re-running the tagger on a copy of the jets (\code{updateJetCollection} with the global-tag AK8PFPuppi
corrections, as in NanoAOD); this copy only provides the 23 \code{globalParT3\_*} columns, matched by \code{srcIdx}, and
does not enter the selection.

\subsection{Generator-level AK8 jets}
\label{sec:genjets}
\code{slimmedGenJetsAK8} jets, with the gen-jet daughters as constituents, are required to have $\pt \geq 180\GeV$;
then only charged constituents are kept, and the jets must have $\pt \geq 200\GeV$, $|\eta| \leq 2.4$ and at least two
charged constituents. No identification, energy correction or smearing is applied. The table \code{GenJet} contains the
same basic columns as \code{RecoJet}, without \code{jerFactor} and the PAT-jet properties.

\subsection{Lepton--jet overlap}
\label{sec:veto}
Jets overlapping with leptons are not removed. Instead, for each channel \code{<c>} listed in
\code{jets.cleaning.channels}, each jet receives a flag \code{veto<c>} that is true if the jet axis lies within
$\Delta R < 0.8$ (\code{cleaningDR}) of a selected lepton of that channel at the jet's level: reco jets are compared
with reco leptons and gen jets with gen leptons. Channel \code{<c>} ignores the jets with \code{veto<c>}; the other
channel is unaffected.

\subsection{Back-to-back requirement}
\label{sec:b2b}
For each channel with \code{backToBack.enabled} (all four channels of \code{ZX\_test}), at each level:
\begin{enumerate}
\item among the jets of that level passing the jet selection and without \code{veto<c>}, the jet with the highest $\pt$
  is taken;
\item $\Delta\phi(\text{jet}, Z) > 2.8$ (\code{dPhiMin}) is required, with $\Delta\phi$ folded into $[0,\pi]$;
\item the $\pt$ asymmetry $A = |\ptj - \ptll| / (\ptj + \ptll) < 0.3$ (\code{asymmetryMax}) is required.
\end{enumerate}
If the leading jet fails these requirements, the next jet is not considered, and \code{B2B<c>} is false. Both cut values
are placeholders. The requirement removes whole events, and only the selected jets are written: per level, the jet
table and the correlators contain the jets selected by the channels that passed, normally one jet, and two only if
both channels pass with different jets. The column \code{b2b<c>} identifies the channel or channels that selected each
jet, and \code{veto<c>} is kept. The requirement is evaluated with the nominal jets only; its re-evaluation for jet
systematic variations is not implemented, and the configuration validation rejects the combination
(Section~\ref{sec:issues}).

%=======================================================================================================================
\section{Energy correlators}
\label{sec:eec}

All correlators are computed for every written jet of both collections, from the selected constituents $i$ (charged,
and PUPPI-weighted for reco jets). The weights use the constituent $\pt^i$ in the laboratory frame, normalised to
$\ptj$, the jet $\pt$ after all corrections. Since $\ptj$ includes the neutral component, the weights do not sum to one.
The correlators are stored unbinned, with one entry per term.

\paragraph{Projected two-point correlator (E2C).} For each pair $\{i, j\}$ with $i \leq j$, including $i = j$, one
entry is stored with
\begin{equation}
  \Delta R_{ij}, \qquad
  w_{ij} = \left(\frac{\pt^i}{\ptj}\right)^{n}\left(\frac{\pt^j}{\ptj}\right)^{n}, \qquad
  s_{ij} = \begin{cases} 2 & i \neq j, \\ 1 & i = j, \end{cases}
\end{equation}
where $\Delta R_{ij} = \sqrt{\Delta\eta^2 + \Delta\phi^2}$ is computed with pseudorapidity and $\Delta\phi$ folded into
$[-\pi,\pi]$, and the symmetry factor $s_{ij}$ is stored separately (\code{symfactor}). The correlator is
$\sum_{i\le j} s_{ij} w_{ij}\,\delta(\Delta R - \Delta R_{ij})$, which equals the sum over ordered pairs with
repetition. Two observables are stored: \code{E2Cn1} ($n = 1$) and \code{E2Cn2} ($n = 2$). The $i = j$ terms have
$\Delta R = 0$ and are removed in the analysis (histogram cut \code{eec.dR > 0}).

\paragraph{One-point correlator (E1C).} For each constituent, one entry is stored with the polar and azimuthal angles
$(\theta_i, \phi_i)$ of its momentum in a frame defined by the jet four-momentum $P$ (after all corrections), the weight
$(\pt^i/\ptj)^n$ with $n = 1$, and symmetry factor 1. If the frame is a rest frame, the constituent is first boosted into
it; $\theta \in [0,\pi]$ and $\phi \in [-\pi, \pi]$ are then measured with respect to the orthonormal axes
$(\hat x, \hat y = \hat z \times \hat x, \hat z)$, with $\hat x$ orthogonalised against $\hat z$. The three frames are:
\begin{description}
\item[\code{E1CLab} (frame \code{jetAxis})] no boost; $\hat z$ along the jet momentum and $\hat x$ in the plane of the
  jet and the beam axis, on the side of the $+z$ beam; $\theta$ is the three-dimensional angle to the jet axis.
\item[\code{E1CHelicity} (frame \code{jetHelicity})] the jet rest frame, reached by a boost along $P$; $\hat z$ along
  the jet direction and $\hat x$ as for \code{jetAxis}, since a boost along the jet direction leaves the jet--beam plane
  unchanged.
\item[\code{E1CCS} (frame \code{jetCS})] the jet rest frame with Collins--Soper axes~\cite{Collins:1977iv}. With unit light-like vectors along
  the two beams, beam $A$ moving along the sign of the jet $p_z$ (the $+z$ beam if $p_z = 0$), and their directions
  $\hat p_A$, $\hat p_B$ in the jet rest frame, $\hat z \propto \hat p_A - \hat p_B$ and
  $\hat x \propto -(\hat p_A + \hat p_B)$.
\end{description}
A jet that is not timelike, or whose axes are degenerate (a jet along the beam), stops the job with an error. The
minimum of two constituents per jet excludes single-particle jets, in whose rest frame the constituent direction is
undefined.

\paragraph{Output.} For each observable and collection (\code{Reco}, \code{Gen}) there are a per-jet table
\code{<Coll><obs>} with \code{E\_jet} ($= \ptj$), \code{n}, \code{N} and \code{nEntries}, and an entries table
\code{<Coll><obs>Entries} with \code{iJet}, the coordinates (\code{dR}, or \code{theta} and \code{phi}), \code{weight}
and \code{symfactor}, all at full floating-point precision. The weights can be renormalised to another jet scale $E'$
with the factor $(E_\mathrm{jet}/E')^{nN}$, where $N = 2$ for E2C and $N = 1$ for E1C. With \code{writeConstituents=1},
the constituents are stored in \code{RecoJetConstituents} and \code{GenJetConstituents} (\code{iJet}, $\pt$, $\eta$,
$\phi$, energy, charge, PDG identifier, the PUPPI weight applied, and \code{srcIdx}).

%=======================================================================================================================
\section{Order of corrections and selections}
\label{sec:order}

CMSSW runs the modules on demand, so the order of the operations is fixed by their data dependencies. This section
states that order for each object, in simulation; in data the generator-level steps and the JER smearing are absent.
Table~\ref{tab:order} lists, for every requirement, the quantity it uses and the corrections already applied to it.

\paragraph{Event-level requirements.} The trigger decision, the primary-vertex requirement and the MET filters use no
corrected object: they are computed from the stored HLT and filter decisions and the vertex collection, independently
of every correction below, and are combined with the dilepton selection only in the channel flag \code{<c>}.

\paragraph{Muons.}
\begin{enumerate}[label=M\arabic*.]
\item The regime is decided from the uncorrected PF $\pt$ ($\geq 200\GeV$: high).
\item The transverse momentum is set to the TuneP $\pt$ (high regime) or kept at the PF $\pt$ (medium regime); the
  direction is that of the PF muon in both cases.
\item Medium regime: if $26 \leq \pt \leq 200\GeV$ (on the $\pt$ before correction), the MuonScaRe scale correction,
  then (simulation) the resolution smearing. High regime: if the TuneP $\pt$ before correction is at least $200\GeV$
  and $|\eta| < 2.4$, the $\kappa_b$ shift, then (simulation) the smearing.
\item Object selection ($\pt > 15\GeV$, $|\eta| < 2.4$) on the corrected muons.
\item Dilepton selection on the corrected muons: lepton $\pt$ thresholds, opposite charge, mass window, choice of the
  pair closest to $m_Z$, then $\ptll \geq 150\GeV$ on the chosen pair.
\end{enumerate}

\paragraph{Electrons.}
\begin{enumerate}[label=E\arabic*.]
\item The MiniAOD four-momentum is used (energy regression of the reconstruction; no scale or smearing correction).
\item Object selection ($\pt > 15\GeV$, $|\eta_\mathrm{SC}| < 2.5$ outside the transition region), then the dilepton
  selection as for muons.
\end{enumerate}

\paragraph{Generator-level leptons.}
\begin{enumerate}[label=G\arabic*.]
\item Rivet particle level: prompt leptons dressed with the photons within $\Delta R < 0.1$.
\item Object selection on the dressed leptons, then the dilepton selection as at reco level.
\end{enumerate}

\paragraph{Reconstructed AK8 jets.}
\begin{enumerate}[label=J\arabic*.]
\item The MiniAOD jet energy correction is removed (raw four-momentum).
\item The V4 L1L2L3Res correction is applied (inputs: area, raw $\pt$, $\eta$, $\rho$).
\item Jets with $\pt < 180\GeV$ after the V4 correction are discarded (before the JER smearing).
\item The constituents are PUPPI-weighted; the jet identification bits are computed.
\item JER smearing of the jet four-momentum, with scale factor and resolution evaluated at the $\pt$ of J2 and the gen
  jet matching done with that $\pt$. The constituents are not changed.
\item Only charged constituents are kept.
\item Jet selection: $\pt \geq 200\GeV$ after J2 and J5, $|\eta| \leq 2.4$, TIGHT identification, at least two charged
  constituents.
\item Overlap flags \code{veto<c>}: $\Delta R$ between the jet axis and the selected leptons of M5 or E2. Neither the
  JEC nor the JER changes the jet direction, and the muon corrections do not change the muon direction.
\item Back-to-back requirement: the leading jet by $\pt$ after J2 and J5; $\Delta\phi$ and the asymmetry use that jet
  $\pt$ and the $\ptll$ of the corrected leptons.
\item Only the selected jets are kept, and the correlators are computed from their constituents after J4 and J6
  (PUPPI-weighted, charged, without JEC or JER), normalised to the jet $\pt$ after J2 and J5; the E1C frames use the jet
  four-momentum after J2 and J5.
\end{enumerate}
The GloParT re-tagging runs on a separate copy of the MiniAOD jets with the global-tag corrections, and does not
feed any step above.

\paragraph{Generator-level AK8 jets.} $\pt \geq 180\GeV$; charged constituents; $\pt \geq 200\GeV$,
$|\eta| \leq 2.4$ and at least two charged constituents; overlap flags against the dressed gen leptons of G2;
back-to-back requirement against the gen $Z$ candidate; correlators as for reco jets.

\paragraph{Event filter and analysis.} The event is written if any \code{B2B<c>} flag is true. In the analysis, the
region's channel flag, then its \code{B2B<c>} flag, then its cuts are applied, all on the stored corrected quantities;
no further correction is applied there.

\begin{table}[htbp]
\centering
\caption{Quantities used by each requirement, and the corrections applied to them beforehand (simulation).}
\label{tab:order}
\small
\begin{tabularx}{\textwidth}{>{\raggedright\arraybackslash}p{4.3cm}>{\raggedright\arraybackslash}p{3.6cm}>{\raggedright\arraybackslash}X}
\toprule
Requirement & Quantity & Corrections applied before \\
\midrule
Trigger, vertex, MET filters & stored decisions, vertex & none \\
Muon regime (M1) & PF $\pt$ & none \\
MuonScaRe range (M3) & PF $\pt$ & none \\
High-$\pt$ correction range (M3) & TuneP $\pt$, $\eta$ & none \\
Muon object and dilepton cuts (M4, M5) & muon $\pt$, $\eta$, $m_{\ell\ell}$, $\ptll$ & momentum assignment, medium- or
  high-regime correction and smearing \\
Electron object and dilepton cuts (E2) & electron $\pt$, $\eta_\mathrm{SC}$, $m_{\ell\ell}$, $\ptll$ & none beyond
  reconstruction \\
Gen lepton and dilepton cuts (G2) & dressed lepton momenta & photon dressing \\
Jet $\pt \geq 180\GeV$ (J3) & jet $\pt$ & V4 JEC \\
JER scale factor, resolution, matching (J5) & jet $\pt$, $\eta$, $\rho$ & V4 JEC \\
Jet selection (J7) & jet $\pt$, $\eta$, ID, constituent count & V4 JEC, JER; PUPPI weights; charged selection \\
Overlap flags (J8) & jet and lepton directions & lepton corrections (directions unchanged) \\
Back-to-back (J9) & jet $\pt$, $\phi$; $\ptll$, $\phi_{\ell\ell}$ & V4 JEC, JER; lepton corrections \\
Correlator weights (J10) & constituent $\pt$; jet $\pt$ & constituents: PUPPI weights only; jet: V4 JEC, JER \\
E1C frames (J10) & jet four-momentum & V4 JEC, JER \\
\bottomrule
\end{tabularx}
\end{table}

%=======================================================================================================================
\section{Event weights, corrections and stored information}
\label{sec:weights}
\label{sec:output}

Table~\ref{tab:corrections} summarises which corrections are applied in the job and which weights are only stored. No
event weight is applied in the job, since \code{ZX\_test} produces no aggregated histograms; all weights are left to
the analysis.

\begin{table}[htbp]
\centering
\caption{Corrections and weights in \code{ZX\_test}.}
\label{tab:corrections}
\small
\begin{tabularx}{\textwidth}{>{\raggedright\arraybackslash}p{4.2cm}>{\raggedright\arraybackslash}X}
\toprule
Item & Treatment \\
\midrule
Generator weight & stored (\code{genWeight} and the Runs-tree sums); used by the analysis \\
LHE, PDF, scale and PS weights & stored (\code{LHEWeight}, \code{LHEPdfWeight}, \code{LHEScaleWeight},
  \code{LHEReweightingWeight}, \code{PSWeight}), with the LHE particles and \code{LHE\_*} quantities as in NanoAOD \\
Pileup weights & stored (\code{Pileup\_weight}, \code{Pileup\_weightUp}, \code{Pileup\_weightDown}, Luminosity POG
  recipe~\cite{LUMpileup}, from the key
  \code{Collisions2022\_359022\_362760\_eraEFG\_GoldenJson} at the in-time true number of interactions); not applied \\
Muon momentum & applied: hybrid assignment, MuonScaRe (medium regime), $\kappa_b$ and smearing (high
  regime)~\cite{MUOcorrections,MUOscare} \\
Muon efficiency corrections & none (\code{muonScaleFactors: \{\}}): no identification or isolation is required \\
Electron energy scale, resolution and efficiency corrections & none; the Run~3 scale (data) and smearing (simulation)
  recommended by the EGamma POG~\cite{EGMsas} are not applied \\
Trigger efficiency corrections & none \\
AK8 jet energy correction & re-applied (V4 L1L2L3Res)~\cite{JMEjerc} \\
AK8 jet energy resolution & nominal hybrid smearing applied in simulation (JRV2)~\cite{JMEjerc} \\
JES and JER uncertainties & not evaluated (no variations) \\
Jet identification & TIGHT required; TIGHTLEPVETO stored~\cite{JMEjetid} \\
AK4 PUPPI jets & stored only (\code{AK4Jet}: V4 correction re-applied, $\pt > 20\GeV$, $|\eta| < 4.7$, with
  \code{jesUncTotal} and flavour-tagging scores) \\
Missing transverse momentum & stored only: PF (\code{MET}, \code{RawMET}), PUPPI (\code{PuppiMET},
  \code{RawPuppiMET}) and \code{GenMET} \\
\bottomrule
\end{tabularx}
\end{table}

Further stored event content: the HLT bits, the primary and secondary vertices (\code{PV}, \code{OtherPV}, \code{SV}),
the generator vertex, six $\rho$ variables, pruned generator particles (\code{GenPart}), and at both levels the
selection flags of Table~\ref{tab:flags}, the $Z$ candidates and the selected leptons (empty unless \code{<c>Sel} passed).
The complete inventory is given in \code{simulation/NANO.md}.

%=======================================================================================================================
\section{Cutflow}
\label{sec:cutflow}

Table~\ref{tab:cutflow} gives the cutflow obtained on one input file per sample, using 30\,000 events or, for the
\code{PTLL}-binned samples, the whole file. The configuration was run with the dilepton $\pt$ and back-to-back
requirements disabled, which writes every event passing a channel at either level together with all its jets, and the
remaining steps were applied with \code{notes/zxtest\_AN/make\_cutflow.py}. On a test file of 1000 events with
$\ptll$ between 400 and 600\GeV, this script reproduces exactly the \code{B2B<c>} flags of the full configuration
(206, 183, 214 and 206 events for reco $\mu\mu$, reco $ee$, gen $\mu\mu$ and gen $ee$). The first step,
\emph{channel flag}, combines the dilepton selection without the $\ptll$ requirement with the trigger, vertex and MET
filters at reco level, and is the dilepton selection alone at gen level. Since events failing it at both levels are not
written, it cannot be divided further with this method. The steps are cumulative; counts are unweighted, and the
percentages are generator-weighted fractions of all processed events.

The fraction of processed events written is 0.30\% for $ZZ\to2\ell2q$ and 0.51\% for $WZ\to2\ell2q$, which corresponds to
roughly $1.5\times10^{5}$ and $7.6\times10^{4}$ events for the full samples of Table~\ref{tab:samples} (unweighted,
extrapolated from one file each), and between 30\% and 48\% for the \code{PTLL}-binned Drell--Yan samples.

\begin{landscape}
\begin{table}[htbp]
\centering
\caption{Cutflow per level and channel on one input file per sample: events (generator-weighted percentage of the
processed events). The last row counts the events passing the full selection in at least one channel at one level,
i.e.\ the events written by \code{ZX\_test}.}
\label{tab:cutflow}
\small
\begin{tabular}{llrrrrrrr}
\toprule
Level, channel & Step & \texttt{ZZto2L2Q} & \texttt{WZto2L2Q} & \texttt{DY\_PTLL200to400} & \texttt{DY\_PTLL400to600} & \texttt{DY\_PTLL600} & \texttt{DY\_incl} & \texttt{TTto2L2Nu} \\
\midrule
\multicolumn{2}{l}{Processed events} & 30000 & 30000 & 10508 & 21385 & 27720 & 30000 & 30000 \\
\midrule
reco Zmumu & channel flag & 2237 (7.48\%) & 2343 (7.82\%) & 2467 (23.5\%) & 5687 (26.6\%) & 7749 (28\%) & 4087 (13.6\%) & 1460 (4.87\%) \\
 & $\pt^{\ell\ell} > 150\GeV$ & 140 (0.469\%) & 210 (0.701\%) & 2425 (23.1\%) & 5579 (26.1\%) & 7546 (27.2\%) & 18 (0.06\%) & 46 (0.148\%) \\
 & $\geq 1$ jet & 69 (0.231\%) & 108 (0.36\%) & 1765 (16.8\%) & 5489 (25.7\%) & 7530 (27.2\%) & 6 (0.02\%) & 13 (0.0437\%) \\
 & $\Delta\phi > 2.8$ & 54 (0.181\%) & 74 (0.247\%) & 1494 (14.2\%) & 4628 (21.6\%) & 6602 (23.8\%) & 4 (0.0133\%) & 8 (0.0269\%) \\
 & asymmetry $< 0.3$ & 54 (0.181\%) & 68 (0.227\%) & 1462 (13.9\%) & 4506 (21.1\%) & 6434 (23.2\%) & 4 (0.0133\%) & 4 (0.0135\%) \\
\midrule
reco Zee & channel flag & 1678 (5.62\%) & 1804 (6.02\%) & 2052 (19.5\%) & 4854 (22.7\%) & 6913 (24.9\%) & 2762 (9.21\%) & 1198 (4\%) \\
 & $\pt^{\ell\ell} > 150\GeV$ & 104 (0.349\%) & 181 (0.604\%) & 2017 (19.2\%) & 4782 (22.4\%) & 6855 (24.7\%) & 20 (0.0668\%) & 52 (0.175\%) \\
 & $\geq 1$ jet & 40 (0.134\%) & 98 (0.327\%) & 1478 (14.1\%) & 4684 (21.9\%) & 6835 (24.7\%) & 5 (0.0167\%) & 26 (0.0874\%) \\
 & $\Delta\phi > 2.8$ & 25 (0.0838\%) & 65 (0.217\%) & 1230 (11.7\%) & 3957 (18.5\%) & 5942 (21.4\%) & 4 (0.0134\%) & 15 (0.0504\%) \\
 & asymmetry $< 0.3$ & 23 (0.0771\%) & 60 (0.2\%) & 1205 (11.5\%) & 3857 (18\%) & 5794 (20.9\%) & 4 (0.0134\%) & 9 (0.0303\%) \\
\midrule
gen Zmumu & channel flag & 2302 (7.69\%) & 2405 (8.02\%) & 2510 (23.9\%) & 5747 (26.9\%) & 7848 (28.3\%) & 4216 (14.1\%) & 1131 (3.78\%) \\
 & $\pt^{\ell\ell} > 150\GeV$ & 143 (0.479\%) & 216 (0.721\%) & 2508 (23.9\%) & 5746 (26.9\%) & 7848 (28.3\%) & 19 (0.0633\%) & 27 (0.0908\%) \\
 & $\geq 1$ jet & 70 (0.235\%) & 108 (0.36\%) & 1877 (17.9\%) & 5669 (26.5\%) & 7839 (28.3\%) & 6 (0.02\%) & 10 (0.0336\%) \\
 & $\Delta\phi > 2.8$ & 56 (0.188\%) & 76 (0.254\%) & 1587 (15.1\%) & 4784 (22.4\%) & 6880 (24.8\%) & 4 (0.0133\%) & 6 (0.0202\%) \\
 & asymmetry $< 0.3$ & 56 (0.188\%) & 72 (0.24\%) & 1559 (14.8\%) & 4686 (21.9\%) & 6762 (24.4\%) & 4 (0.0133\%) & 2 (0.00673\%) \\
\midrule
gen Zee & channel flag & 2095 (7.01\%) & 2202 (7.35\%) & 2255 (21.5\%) & 5230 (24.5\%) & 7421 (26.8\%) & 3391 (11.3\%) & 960 (3.22\%) \\
 & $\pt^{\ell\ell} > 150\GeV$ & 127 (0.426\%) & 211 (0.704\%) & 2253 (21.4\%) & 5229 (24.5\%) & 7419 (26.8\%) & 20 (0.0668\%) & 19 (0.0639\%) \\
 & $\geq 1$ jet & 46 (0.154\%) & 108 (0.36\%) & 1718 (16.3\%) & 5130 (24\%) & 7410 (26.7\%) & 6 (0.02\%) & 10 (0.0336\%) \\
 & $\Delta\phi > 2.8$ & 29 (0.0972\%) & 70 (0.234\%) & 1445 (13.8\%) & 4346 (20.3\%) & 6451 (23.3\%) & 4 (0.0134\%) & 7 (0.0235\%) \\
 & asymmetry $< 0.3$ & 27 (0.0905\%) & 67 (0.224\%) & 1421 (13.5\%) & 4256 (19.9\%) & 6333 (22.8\%) & 4 (0.0134\%) & 5 (0.0168\%) \\
\midrule
\multicolumn{2}{l}{Written (any level, channel)} & 89 (0.298\%) & 152 (0.507\%) & 3140 (29.9\%) & 9137 (42.7\%) & 13319 (48\%) & 8 (0.0267\%) & 16 (0.0538\%) \\
\bottomrule
\end{tabular}

\end{table}
\end{landscape}

%=======================================================================================================================
\section{Analysis regions}
\label{sec:regions}

The columnar analysis reads the merged NANO with the configuration \code{analysis/configs/zxtest.yaml}. A region is one
channel at one level, optionally narrowed by further cuts, and is built as follows.
\begin{enumerate}
\item Events with \code{Selection\_<c>} (\code{GenSelection\_<c>} at gen level) are kept (cutflow step
  \code{channel}).
\item If the NANO contains the back-to-back flag \code{B2B<c>}, events without it are removed (cutflow step
  \code{backToBack}). This removes events written because another channel or level passed, which have no jet of
  channel \code{<c>}.
\item The objects of the region are the channel's $Z$ candidate (\code{boson}), its selected leptons (\code{leptons}),
  the missing transverse momentum (PUPPI at reco, generator-level at gen), the jets of the channel (\code{jets}), and
  their correlator entries (\code{eec}). The jets of the channel are those with \code{b2b<c>} if the NANO contains
  that column, and otherwise those without \code{veto<c>}. NANO files without the back-to-back columns (such as the
  first production) are therefore processed as before.
\item The region's cuts are applied in order. Each cut is an expression on these objects with a level fixed by its
  form: an event-level cut removes events, a jet-level cut removes jets together with their correlator entries, and a
  lepton-level cut removes the event if any selected lepton fails. Cuts on individual correlator entries are allowed in
  histogram definitions only.
\end{enumerate}
The event weight is \code{genWeight}, multiplied by $\sigma L/\code{genEventSumw}$ once cross sections are available;
correlator entries are filled with the event weight times \code{weight} times \code{symfactor}. On the test file of
Section~\ref{sec:cutflow}, the \code{backToBack} step reproduces the \code{B2B<c>} counts exactly and removes 2 or 3
events per region.

\begin{table}[htbp]
\centering
\caption{Regions of \code{analysis/configs/zxtest.yaml}, in addition to the channel and back-to-back steps.}
\label{tab:regions}
\small
\begin{tabularx}{\textwidth}{lll>{\raggedright\arraybackslash}X}
\toprule
Region & Level & Channel & Additional cuts, in order \\
\midrule
\code{ZmumuReco} & reco & \code{Zmumu} & none \\
\code{ZeeReco} & reco & \code{Zee} & none \\
\code{ZmumuGen} & gen & \code{Zmumu} & none \\
\code{ZeeGen} & gen & \code{Zee} & none \\
\code{ZmumuBoostedReco} & reco & \code{Zmumu} & both muons pass the medium identification (lepton level);
  $\ptll > 200\GeV$ (event level); jet $\pt > 300\GeV$, $|\eta^\mathrm{jet}| < 2.0$ and $\Delta R(\text{jet}, Z) > 2.0$
  (jet level); at least one jet (event level) \\
\bottomrule
\end{tabularx}
\end{table}

The histograms are: $\ptll$, $m_{\ell\ell}$, the number of jets, lepton $\pt$, jet $\pt$ and $\eta$, and, at reco level
only, the jet soft-drop mass, $\tau_{21} = \tau_2/\tau_1$ and the GloParT $W$-versus-QCD score; E2C ($n = 1, 2$) versus
$\Delta R$ (50 logarithmic bins between $10^{-3}$ and 1, without the $\Delta R = 0$ entries) and jet $\pt$ (bin edges
200, 300, 400, 600, 1000 and 3000\GeV); and E1C versus $(\theta, \phi)$ in the three frames ($\theta$ in 40 logarithmic
bins between $10^{-3}$ and 1 for the laboratory frame and in 40 bins between 0 and $\pi$ for the rest frames; $\phi$ in
40 bins between $-\pi$ and $\pi$).

%=======================================================================================================================
\section{Validation}
\label{sec:validation}

The following checks were made on the implementation; details are recorded in \code{simulation/JOURNAL.md}.
\begin{itemize}
\item Unit tests of all modules and of the framework-independent selection, jet and correlator code
  (\code{scram b runtests}), and of the columnar analysis; deliberately altered versions of the logic are verified to
  fail them.
\item Correlators: every unbinned entry (reco and gen, five observables, four samples) agrees with a NumPy
  recomputation from the stored constituents (\code{configs/checkUnbinned.py}), with maximum differences of
  $9\times10^{-5}$ in the coordinates and $3\times10^{-7}$ in the relative weights.
\item Re-applied AK8 jet energy corrections: no mismatch in raw $\pt$ and energy fractions with respect to an
  FWLite/correctionlib computation, on 802 simulated and 428 data jets. GloParT scores identical to those of the standard
  NanoAOD sequence (1056 jets, 23 columns).
\item Dilepton $\pt$ requirement: on 2000 inclusive Drell--Yan events, 8 events are written instead of 513; they are
  exactly the events with a reco or gen $Z$ candidate above $150\GeV$ in the run without the requirement.
\item Back-to-back requirement (\code{JetSetup/test/validation/checkBackToBack.py}): a run with the requirement is
  compared with an independent recomputation from a run without it. On 1000 events with $\ptll$ between 400 and
  600\GeV, 431 of 517 events are kept; the selected jets (809, all columns), their correlator entries (4045
  jet--observable pairs) and the flags agree exactly, and the other 903 columns are unchanged. With the requirement
  disabled, the configuration builds a CMSSW process identical to the one before the feature was added.
\item Columnar analysis: histogram integrals agree with an independent uproot computation on NANO files with and without
  the back-to-back columns.
\end{itemize}

%=======================================================================================================================
\section{Known issues and open items}
\label{sec:issues}

\begin{enumerate}
\item \textbf{AK8 jets without daughters.} In 3 of the 232 $t\bar t$ jobs of the first production, a
  \code{slimmedJetsAK8} jet had a MiniAOD $\pt$ below $170\GeV$, so that its daughters were not stored, but a raw $\pt$
  of $326.8\GeV$ (in both events inspected), which the re-applied V4 correction raises to about $348\GeV$ (in all three
  failures), above the $180\GeV$ threshold; the job then stops with an error. The identical raw $\pt$ in independent
  events suggests a MiniAOD artefact. These jobs were excluded from the merge; the cause and the treatment of such jets
  are open, and the same failure can occur in the next production.
\item \textbf{Back-to-back requirement with jet variations.} Not implemented; it is needed before jet systematic
  variations are enabled. The required steps are listed in \code{simulation/backToBack\_plan.md}.
\item \textbf{Muon corrections.} The formula and sign of the $\kappa_b$ correction, the smearing and its systematic
  variations, the scale correction above $200\GeV$ and the definition of the resolution variation are provisional,
  pending the MUO POG (\code{simulation/EventSelection/README.md}).
\item \textbf{Lepton identification, isolation, trigger and corrections.} By design, none is applied at reco level;
  all lepton and trigger efficiency corrections are still missing. The Run~3 electron energy scale (data) and smearing
  (simulation) recommended by the EGamma POG~\cite{EGMsas} are not applied. The muon scale and resolution corrections
  are applied to muons without identification or isolation, although they were derived for tightly identified and
  isolated muons~\cite{MUOscare}.
\item \textbf{Electron no-isolation MVA identification.} The \code{mvaNoIsoWP80} and \code{mvaNoIsoWP90} columns are
  read from \code{slimmedElectrons}, whose Run~3 no-isolation MVA decisions include isolation in the
  barrel~\cite{EGMmvaid}; the POG fix (re-running the identification on MiniAOD) is not applied, so these columns must
  not be used as no-isolation identifications.
\item \textbf{Cross sections.} Not yet collected; they are needed for the luminosity normalisation and for combining
  the \code{PTLL}-binned Drell--Yan samples with the inclusive one, which also requires a generator-level overlap
  removal.
\item \textbf{Sample content.} Whether \code{ZZto2L2Q} and \code{WZto2L2Q} include $Z\to\tau\tau$ decays, and the
  definition of \code{PTLL}, are to be confirmed from the generator configuration (Section~\ref{sec:samples}).
\item \textbf{Placeholder values.} $\ptll \geq 150\GeV$, $\Delta\phi > 2.8$ and $A < 0.3$ are development choices.
\end{enumerate}

\printbibliography[heading=bibintoc]

\end{document}
